\documentclass[17pt,a4paper,landscape]{extarticle}
\usepackage[margin=2.35cm,top=0.12cm,left=1.05cm]{geometry}
\usepackage{amsmath}
\usepackage{amssymb}
\usepackage{tasks}
\usepackage{tikz}
\usepackage{xcolor}
\usepackage[sfdefault,lf]{FiraSans}
\renewcommand{\familydefault}{\sfdefault}
\setlength{\parindent}{0pt}
\pagestyle{empty}
\settasks{label=\textbf{\arabic*.}, label-width=2.2em, item-indent=2.95em, column-sep=2.1cm, after-item-skip=4.7em}
\renewcommand{\arraystretch}{1.15}
\everymath{\displaystyle}

\begin{document}
\sffamily
\boldmath

\boldmath

\boldmath

\boldmath

{\LARGE \textbf{Area of trapezium — Excellence}}
\begin{tasks}(3)
\task {Identify and circle the two parallel sides in the drawn trapezium.}
\task {Calculate the area of a trapezium with parallel sides of 16 m and 3 m, and a perpendicular height of 7 m.}
\task {Explain why the area formula for a trapezium can be thought of as "the average of the parallel sides multiplied by the height."}
\task {A shape consists of a trapezium with parallel sides 7 m and 5 m, height 18 m, attached to a rectangle of length 12 m and width 3 m. Find the total area.}
\task {The parallel sides of a trapezium are \mbox{$x$} m and \mbox{$2x$} m, where \mbox{$x = 20$}. The perpendicular height is \mbox{$x + 2$} m. Calculate the area.}
\task {Identify and circle the two parallel sides in the drawn trapezium.}
\task {Calculate the area of a trapezium with parallel sides of 6 m and 15 m, and a perpendicular height of 4 m.}
\task {Explain why the area formula for a trapezium can be thought of as "the average of the parallel sides multiplied by the height."}
\task {A shape consists of a trapezium with parallel sides 7 m and 20 m, height 16 m, attached to a rectangle of length 11 m and width 7 m. Find the total area.}
\task {The parallel sides of a trapezium are \mbox{$x$} m and \mbox{$2x$} m, where \mbox{$x = 6$}. The perpendicular height is \mbox{$x + 2$} m. Calculate the area.}
\task {Identify and circle the two parallel sides in the drawn trapezium.}
\task {Calculate the area of a trapezium with parallel sides of 12 m and 1 m, and a perpendicular height of 19 m.}
\task {Explain why the area formula for a trapezium can be thought of as "the average of the parallel sides multiplied by the height."}
\task {A shape consists of a trapezium with parallel sides 5 m and 16 m, height 8 m, attached to a rectangle of length 12 m and width 5 m. Find the total area.}
\task {The parallel sides of a trapezium are \mbox{$x$} m and \mbox{$2x$} m, where \mbox{$x = 15$}. The perpendicular height is \mbox{$x + 2$} m. Calculate the area.}
\task {Identify and circle the two parallel sides in the drawn trapezium.}
\end{tasks}

\end{document}
